\subsection{LLM-based code generation and program synthesis}
\label{subsec:rw0}

The automatic production of executable code from a higher-level description is a long-standing goal of computer science, and the recent confluence of classical program synthesis with large language models (LLMs) has dramatically reshaped its landscape. Understanding where this article's reference pipeline sits requires distinguishing two intellectual lineages that the pipeline deliberately re-couples: \emph{specification-driven} synthesis, in which a precise formal artefact defines the target, and \emph{description-driven} generation, in which an LLM maps natural-language (NL) intent to plausible code.

The synthesis lineage is the older of the two. Deductive synthesis derives a program from a constructive proof that a formal specification is satisfiable~\cite{manna1980deductive}, while inductive and example-driven approaches infer programs from input--output behaviour, as in programming-by-example systems that synthesize string transformations from a handful of demonstrations~\cite{gulwani2011flashfill}. Counterexample-guided inductive synthesis (CEGIS) and the closely related \emph{sketching} paradigm partition the problem into a synthesizer that proposes a candidate and a verifier that either accepts it or returns a counterexample to refine the search~\cite{solarlezama2008sketching}. The defining property of this lineage is that the target is a machine-checkable contract and the loop is \emph{closed}: a candidate is admitted only when an oracle certifies it. That oracle, however, has historically been the bottleneck---formal specifications are expensive to author, and synthesis scales poorly beyond small, well-circumscribed problems.

The description-driven lineage emerged with the application of transformer language models to source code. Codex demonstrated that a model trained on public repositories could generate functionally correct programs from docstrings, and introduced the HumanEval benchmark together with the now-standard \emph{pass@k} metric for functional correctness~\cite{chen2021codex}; the contemporaneous MBPP benchmark established a complementary set of entry-level Python tasks~\cite{austin2021mbpp}. Competition-level results followed, with AlphaCode reaching median human performance on programming contests by massive sampling and filtering~\cite{li2022alphacode}. These systems inverted the economics of the synthesis lineage---NL intent is cheap and abundant---but at the cost of abandoning the closed loop: correctness is estimated against a finite, model-adjacent test suite rather than proved against a contract. Subsequent scrutiny exposed the fragility of this estimate. Benchmark test suites are frequently too weak to detect plausible-but-incorrect solutions, and strengthening them materially lowers reported pass rates~\cite{liu2023evalplus}. More fundamentally, LLM code generation is prone to \emph{hallucination}: confidently produced fabrications such as non-existent APIs, intent conflicts, and inconsistent context dependencies, which evade superficial review precisely because the surrounding code is well-formed~\cite{zhang2025hallucination}. Because the generator and its informal oracle share a model and a training distribution, the verification is effectively self-grading.

The frontier of the description-driven lineage has moved from function-level snippets to repository-level, \emph{agentic} software engineering. SWE-bench reframed evaluation around resolving real GitHub issues against full codebases with project test suites as the oracle~\cite{jimenez2024swebench}, and a wave of agents followed: SWE-agent showed that a purpose-built agent--computer interface substantially improves an LLM's ability to navigate and edit a repository~\cite{yang2024sweagent}, while AutoCodeRover combined LLM reasoning with structure-aware code search and spectrum-based fault localization to localize and patch issues~\cite{zhang2024autocoderover}. These agents are impressive, but they preserve and indeed amplify the two structural defects of the lineage: the specification of intent remains the NL issue text, and the acceptance gate remains an executable test suite---typically the project's own tests, which the agent may also be permitted to consult or modify, reintroducing circularity at the repository scale.

The gap this article addresses lies precisely at the seam between these lineages. Agentic generation has solved the authoring-cost problem that crippled classical synthesis, but in doing so it discarded the machine-checkable contract and the closed, counterexample-driven loop that made synthesis trustworthy. This article's contribution is to restore that loop in an LLM-native setting by interposing a \emph{formal-specification contract layer} that serves the dual role of input intent and verification oracle, with the contract verified against \emph{all} inputs by OpenJML extended static checking rather than sampled by tests. Crucially, the oracle is not the generator: an independent-authority constraint and provenance- and assurance-tagged obligations make the ``verified'' claim auditable, directly answering the self-grading critique that the reliability literature raises~\cite{liu2023evalplus,zhang2025hallucination}. The synthesis--verification loop is CEGIS-shaped~\cite{solarlezama2008sketching}, but operates at the granularity of agentic edits~\cite{yang2024sweagent,zhang2024autocoderover} and degrades gracefully when verification does not terminate. Finally, where the agentic benchmarks are greenfield-issue-centric~\cite{jimenez2024swebench}, this article targets the brownfield reality of most software through multi-source intent triangulation---reconciling code, tickets, and documentation into a provenance-tagged contract---and lifts the author's prior compositional inference~\cite{edmonds_article3} into a lifecycle capability, so that a callee-spec change recomputes caller obligations without re-running them.
